%% =====================================================================
%%  Benchmarking General and Medical Small Language Models for
%%  Spanish Obstetric Question Answering: The Impact of Fine-Tuning
%%
%%  IEEE conference template (IEEEtran.cls).
%%  Build: pdflatex -> bibtex -> pdflatex -> pdflatex
%%
%%  Class options:
%%    conference : IEEE conference/symposium style (two columns, full-width
%%                 title block, "Abstract" run-in head, "Index Terms").
%%    Use [journal] instead for IEEE Transactions/Journals submission.
%%
%%  Deliberately NOT loaded (incompatible with IEEEtran or forbidden by
%%  the IEEE template): geometry, caption, float, placeins, newtxtext,
%%  newtxmath, inputenc, enumitem, dblfloatfix, stfloats, xcolor.
%%
%%  stfloats is omitted on purpose: the IEEEtran HOWTO documents it as a
%%  "very invasive package" needed only for double-column floats at the
%%  BOTTOM of a page (\begin{figure*}[!b]). Every full-width float here is
%%  placed at the top, so it would add output-routine risk with no benefit.
%% =====================================================================
\documentclass[conference]{IEEEtran}

% --- Fonts and encoding ---------------------------------------------
\usepackage[T1]{fontenc}

% --- Mathematics ----------------------------------------------------
\usepackage{amsmath,amssymb,amsfonts}

% --- Graphics, floats and tables ------------------------------------
\usepackage{graphicx}
% needspace keeps a short paragraph from being split by a full-width float:
% it breaks the column instead of letting the text straddle the float.
\usepackage{needspace}
\usepackage{booktabs}          % \toprule \midrule \bottomrule \cmidrule

\usepackage{array}
\usepackage[caption=false]{subfig}  % sub-figures; caption=false is required
                                    % because the caption package must not be
                                    % used together with IEEEtran
\usepackage{threeparttable}    % professional table notes
\usepackage{textcomp}

% Avoid stacking multiple figures in one column and keep the float sequence
% close to the relevant explanations.
\setcounter{topnumber}{1}
\setcounter{dbltopnumber}{2}

% --- Citations and hyperlinks ---------------------------------------
\usepackage{cite}              % IEEE-recommended numeric citation handling

% --- Typographic fine-tuning ----------------------------------------
% microtype is loaded before hyperref so that both patch the output routine
% and the PDF-string machinery in the documented order.
\usepackage{microtype}

\usepackage[hidelinks]{hyperref}

% --- Project-local shorthands ---------------------------------------
% \code marks technical identifiers (model ids, paths, configuration keys).
% It is declared as a URL-style command instead of plain \texttt so that a
% long monospaced token may break instead of overflowing its column.
% Two consequences to keep in mind:
%   * the argument is read verbatim, so write the underscore directly;
%   * it is URL-like, so spaces inside the argument are NOT preserved. Use it
%     for single tokens only.
% \UrlBreaks below is url's default break list with two edits: the underscore
% is removed (breaking an identifier at "_" reads as two separate words,
% e.g. "clinical_" / "score") and the hyphen is added (url does not break on
% hyphens by default, which matters for long model ids).
\def\UrlBreaks{\do\.\do\@\do\\\do\/\do\!\do\|\do\;\do\>\do\]%
  \do\)\do\,\do\?\do\'\do+\do\=\do\#\do\-}
\DeclareUrlCommand\code{\urlstyle{tt}}
\newcommand{\dash}{---}

\begin{document}

% =====================================================================
%  Title, authors and affiliations
%  IEEE rule: title in upper and lower case, never all upper case, and no
%  symbols, math, footnotes or citations inside the title.
%
%  Author block: each author gets its own \IEEEauthorblockN plus one
%  \IEEEauthorblockA, separated by \and. IEEEtran lays these out in columns,
%  matching the official IEEE conference template (name, then dept./org. in
%  italic, then City, Country upright, then email address or ORCID).
% =====================================================================
\title{Benchmarking General and Medical Small Language Models for
Spanish Obstetric Question Answering: The Impact of Fine-Tuning}

\author{%
\IEEEauthorblockN{Jhon H. Mej\'ia-Muñoz}
\IEEEauthorblockA{\textit{Facultad de Ingenier\'ia, IUE}\\
Envigado, Colombia\\
0009-0008-7768-4072}
\and
\IEEEauthorblockN{Nicol\'as Hoyos-Giraldo}
\IEEEauthorblockA{\textit{Facultad de Ingenier\'ia, IUE}\\
Envigado, Colombia\\
0009-0008-7768-4072}
\and
\IEEEauthorblockN{Daniel Betancur-Vasquez}
\IEEEauthorblockA{\textit{Facultad de Ingenier\'ia, IUE}\\
Envigado, Colombia\\
0000-0003-4834-7268}
\and
\IEEEauthorblockN{Kevin D. Trujillo}
\IEEEauthorblockA{\textit{Facultad de Ingenier\'ia, IUE}\\
Envigado, Colombia\\
0009-0001-7678-2681}
}

\maketitle

% =====================================================================
%  Abstract and index terms
%  IEEE order: \maketitle, then abstract, then IEEEkeywords.
% =====================================================================
\begin{abstract}
Small Language Models (SLMs) offer a computationally efficient alternative
to large language models for clinical question answering; however, evidence
on domain adaptation for Spanish obstetric applications remains limited.
This study compares the general-purpose Gemma~4 model and the medically
pretrained MedGemma~1.5 before and after grounded QLoRA adaptation using
MaternaQA-es, while examining No RAG, Hybrid retrieval, and
independent HyDE as experimental factors. All 24 scenario-condition runs
were evaluated across two 10-question scenarios, yielding 240 final
condition-item outputs. Within this exploratory evaluation, on MaternaQA-es,
Gemma~4 Base with Hybrid retrieval achieved the highest observed Answer Correctness (0.649),
while Gemma~4 QLoRA with No RAG and Hybrid retrieval achieved the highest
Answer Relevancy (0.776) and Semantic Similarity (0.838), respectively.
On the obstetric challenge set, Gemma~4 Base with Hybrid retrieval achieved
the highest Answer Correctness (0.457), MedGemma~1.5 QLoRA with No RAG the
highest Answer Relevancy (0.913), and Gemma~4 Base and Gemma~4 QLoRA with HyDE tied
in Semantic Similarity (0.707). HyDE achieved higher RAGAS Context Precision
than Hybrid in seven of eight model-scenario comparisons; effects on answer
quality were heterogeneous and model-dependent. Under No RAG, QLoRA increased
Answer Relevancy in all four model-scenario comparisons, while changes in
Answer Correctness and Semantic Similarity were mixed. Given the limited
evaluation of 10 questions per scenario, these findings should be
interpreted as preliminary descriptive patterns.
\end{abstract}

\begin{IEEEkeywords}
Small language models, QLoRA, retrieval-augmented generation, HyDE,
obstetric question answering, Spanish clinical NLP, model benchmarking.
\end{IEEEkeywords}

% =====================================================================
\section{Introduction}
% =====================================================================
The rapid adoption of Large Language Models (LLMs) has transformed the
development of clinical decision-support systems and medical
question-answering (QA) applications. Nevertheless, their computational
requirements limit deployment in many healthcare environments, motivating
increasing interest in Small Language Models (SLMs), which offer lower
inference costs, reduced hardware requirements, and easier local
deployment while maintaining competitive performance in domain-specific
tasks~\cite{pingua2025medical,gemma2024}. In medical applications,
supervised fine-tuning has emerged as a practical strategy for adapting
compact language models to specialized clinical domains while preserving
computational efficiency.

Recent studies have shown that domain-specific fine-tuning and retrieval
can improve medical question answering under resource
constraints~\cite{bora2024systematic}. Work on fine-tuned clinical
embedding models likewise demonstrates the value of domain-specific
retrieval, but remains outside Spanish obstetric QA~\cite{arzideh2026improving}.
Recent obstetric applications include simulated-case decision support,
emergency-response evaluation, and neonatal growth
prediction~\cite{psilopatis2026role,arslan2025evaluation,xu2026obstet}.
Reviews of maternal-health RAG emphasize the lack of standardized
benchmarks and clinically validated evaluation
protocols~\cite{noguera2025applications}. This gap is especially relevant
in obstetric care, where inaccurate recommendations may affect maternal
and fetal safety~\cite{pacheco2023inteligencia}.

An important question therefore remains: does medical pretraining continue
to provide measurable advantages after domain-specific fine-tuning?
Medically pretrained SLMs should encode biomedical terminology, clinical
concepts, and reasoning patterns. Yet, after both a general-purpose and a
medical model receive identical specialized supervision, it remains unclear
whether prior medical knowledge provides an additional benefit or whether
the general-purpose model can achieve comparable or superior performance.

To investigate this question, this study uses MaternaQA-es, a Spanish
obstetric QA resource containing 5,727 synthetic question-answer pairs
generated from 63 Spanish-language clinical documents. The dataset is
publicly available through
Hugging Face~\cite{maternaqa2026}. The benchmark compares Gemma~4 and
MedGemma~1.5 before and after identical grounded QLoRA adaptation.

This work makes three contributions. First, it compares general-purpose and
medically pretrained SLMs for Spanish obstetric QA. Second, it evaluates
the effect of domain-specific QLoRA adaptation under matched training
conditions. Third, it provides a factorial benchmark incorporating
retrieval-augmented inference and item-level evaluation.

% =====================================================================
\section{Methodology}
% =====================================================================
\subsection{Study Design and Research Questions}
The study uses a controlled, offline factorial benchmark. Model family has
two levels: Google Gemma~4 E2B instruction-tuned
(\code{google/gemma-4-E2B-it}) and Google MedGemma~1.5 4B instruction-tuned
(\code{google/medgemma-1.5-4b-it}). Parameter adaptation has two levels:
base weights without an adapter and grounded QLoRA. Retrieval strategy has
three planned levels: the No RAG control, canonical
Hybrid retrieval, and canonical independent HyDE. Their Cartesian product
defines 12 planned conditions. Throughout the paper, ``Base'' means no
QLoRA adapter, and ``No RAG'' means no retrieved context (the model answers
from the question alone); the two are independent and can be combined.

The unit of analysis is an item-level paired \code{qa_id}. The primary
QLoRA estimand is the paired difference between QLoRA and base outputs
within a fixed model family and retrieval strategy. Retrieval contrasts
compare strategies within a fixed family and adaptation state. The
interaction is the QLoRA-minus-base difference under Hybrid minus that
under No RAG, within a model family; analogous contrasts compare HyDE
with No RAG and the two model families.

\subsection{Dataset, Provenance, and Split Integrity}
MaternaQA-es contains 5,727 Spanish obstetric question-answer pairs derived
from 63 clinical documents. Its document-level partitions contain 52
training PDFs (5,093 items), 2 validation PDFs (306 items), and 3 held-out
test PDFs (328 items). The exploratory evaluation subset comprised the
first 10 records of the held-out test file used in the experiment, rather
than the complete test partition. Test sources do not overlap training or
validation sources, and \code{qa_id} is retained as the pairing
key~\cite{maternaqa2026}.

The primary grounded adaptation uses the \code{sft_grounded}
representation, in which the source context associated with an item is
concatenated with the question.

The same 10 held-out test questions were processed across all model and
retrieval conditions. The 2,268-chunk retrieval corpus includes the source
material associated with these questions. Therefore, the RAG evaluation
measures the system's ability to retrieve and use relevant evidence when
that evidence is available in the corpus, rather than strict generalization
to completely unseen source documents. This inference-time access is
distinct from leakage into model parameters, because the corresponding test
documents were not used during fine-tuning.

\subsection{Protocol Evaluation Scenarios}
Two scenarios each comprise 10 questions evaluated across all 12 conditions:
the obstetric challenge set (\code{datasets/sample10.jsonl}) and the
processed held-out MaternaQA-es test subset. The latter retains QA metadata
and source-reference information; the challenge set provides no gold
reference-chunk identifiers. MaternaQA-es is publicly available on
Hugging Face~\cite{maternaqa2026}.

In both scenarios, the answer prompt contains only the question before
retrieval. Hybrid and HyDE add retrieved corpus contexts, whereas No RAG
receives none. Both scenarios support reference-answer metrics; only
MaternaQA-es supports reference-chunk identity checks.

Therefore, the two scenarios are treated as complementary evaluation
settings rather than directly comparable samples from the same input
distribution. Absolute metric differences between MaternaQA-es and the
obstetric challenge set are not interpreted as effects of fine-tuning or
retrieval, because scenario membership is confounded with differences in
input construction and source-context availability. Model and retrieval
comparisons are consequently emphasized within each scenario.

\subsection{Model Families and Experimental Conditions}
The benchmark compares Google Gemma~4 E2B instruction-tuned with Google
MedGemma~1.5 4B instruction-tuned. Base conditions load the published
model without an adapter. QLoRA conditions load
\code{iue-edu/MaternaCare-ES-gemma4-qlora} or
\code{iue-edu/MaternaCare-ES-medgemma-qlora} on the corresponding base
model. The answer model, adaptation state, and test items are held fixed across
retrieval-strategy comparisons within each model family; the retrieval
strategy determines whether retrieved evidence is supplied.

Table~\ref{tab:factorial} lists the complete prespecified design. The final
result artifacts covered No RAG, Hybrid, and HyDE for all four
model/adaptation combinations in both scenarios; therefore, all 24
scenario-condition runs were represented in the analysis.

% --- Table I: single column. The three factor columns already encode the
% condition, so the old fifth column restated each row in prose and forced a
% full-width float that cost vertical space on every page it touched.
\begin{table}[t]
\footnotesize
\centering
\begin{threeparttable}
\caption{Prespecified Factorial Experiment Matrix}
\label{tab:factorial}
\begin{tabular*}{\columnwidth}{@{\extracolsep{\fill}}llll@{}}
\toprule
\textbf{ID} & \textbf{Model family} & \textbf{Adaptation} & \textbf{Retrieval} \\
\midrule
C01 & Gemma 4     & Base  & No RAG \\
C02 & Gemma 4     & Base  & Hybrid \\
C03 & Gemma 4     & Base  & HyDE \\
C04 & Gemma 4     & QLoRA & No RAG \\
C05 & Gemma 4     & QLoRA & Hybrid \\
C06 & Gemma 4     & QLoRA & HyDE \\
C07 & MedGemma 1.5 & Base  & No RAG \\
C08 & MedGemma 1.5 & Base  & Hybrid \\
C09 & MedGemma 1.5 & Base  & HyDE \\
C10 & MedGemma 1.5 & QLoRA & No RAG \\
C11 & MedGemma 1.5 & QLoRA & Hybrid \\
C12 & MedGemma 1.5 & QLoRA & HyDE \\
\bottomrule
\end{tabular*}
\begin{tablenotes}[flushleft]
\footnotesize
\item[] \textit{Note.} \textit{Base} denotes published weights without an
adapter; \textit{QLoRA} denotes the grounded adapter; \textit{No RAG} denotes
no retrieved context.
\end{tablenotes}
\end{threeparttable}
\end{table}

\subsection{Grounded QLoRA Fine-Tuning}
\label{sec:qlora}
Both model families used grounded QLoRA~\cite{dettmers2023qlora} with
4-bit base weights and the shared settings in
Table~\ref{tab:qlora-config}. Training and validation partitions, grounded
inputs, and optimization settings were matched across model families.

% --- Table II: single column; only the base model differs -----------------
\begin{table}[t]
\footnotesize
\centering
\begin{threeparttable}
\caption{Grounded QLoRA Configuration Used for Both Model Families}
\label{tab:qlora-config}
\begin{tabular}{@{}>{\raggedright\arraybackslash}p{\dimexpr0.43\columnwidth-\tabcolsep\relax}%
                  >{\raggedright\arraybackslash}p{\dimexpr0.57\columnwidth-\tabcolsep\relax}@{}}
\toprule
\textbf{Parameter} & \textbf{Shared setting} \\
\midrule
Base model & Gemma 4 E2B: \code{google/gemma-4-E2B-it}\newline
             MedGemma 1.5 4B: \code{google/medgemma-1.5-4b-it} \\
Training representation & sft\_grounded \\
Train / validation examples & 5,093 / 306 \\
Quantization & 4-bit NF4, double quantization \\
LoRA rank ($r$) & 16 \\
LoRA alpha & 16 \\
LoRA dropout & 0.05 \\
Target modules & All linear (all-linear) \\
Learning rate & $2 \times 10^{-4}$ \\
Batch size & 1 per device \\
Gradient accumulation steps & 8 \\
Effective batch size & 8 \\
Epochs / training steps & 2 / 1,274 \\
Optimizer & AdamW 8-bit \\
Learning-rate scheduler & Cosine \\
Maximum sequence length & 1,024 tokens \\
Packing & Disabled \\
Random seed & 3407 \\
\bottomrule
\end{tabular}
\end{threeparttable}
\end{table}

The \code{sft_grounded} examples use the conversational
\code{Contextofuente} representation: source context is pre-associated
with each training question, unlike inference-time RAG selection from the
2,268-chunk corpus. Only QLoRA conditions load an adapter.

\subsection{Retrieval-Augmented Generation Protocol}
\label{sec:rag-protocol}
Retrieval strategy~\cite{lewis2020retrieval} controls evidence preparation:
No RAG supplies none, Hybrid fuses dense and sparse retrieval, and
independent HyDE uses hypothetical text only for query expansion.

\subsubsection{Corpus Provenance and Chunk Construction}
The retrieval corpus contains 2,268 chunks from curated obstetric documents.
Its corpus-filtering \code{clinical_score} adds one point per clinical term
found as a substring of normalized text (two for multiword terms) and one
per action term. Cumulative bonuses add two points at each Unicode word-count
threshold of 120 and 350; a reference-line ratio of at least 0.35 subtracts
three points, with the final score floored at zero. The ratio counts
reference-matching nonempty lines over all nonempty lines. Normalization
applies Unicode normalization, removes combining marks, lowercases text,
replaces non-ASCII letters and punctuation with spaces, and collapses whitespace.

\subsubsection{Dense and Sparse Retrieval}
Dense retrieval uses \code{BAAI/bge-m3}~\cite{chen2024bge} with
SentenceTransformer-normalized vectors on CPU (\code{batch_size=32}) in a
persistent local LanceDB table~\cite{lancedb2026vector}. Search is exhaustive
exact cosine search, without an approximate index; default depth is five.
Sparse retrieval uses Okapi BM25 (\code{rank_bm25/BM25Okapi}) over Unicode
lowercase word tokens. Outputs are recorded before answer generation for
separate retrieval-quality and latency analysis.

\subsubsection{Canonical Hybrid Retrieval}
Hybrid retrieval fuses dense and BM25 rankings by reciprocal rank fusion
with \code{rank_constant=60}~\cite{cormack2009reciprocal}. Candidate depth
is the greater of four times the requested final depth and 50, capped by
corpus size. Stable chunk IDs break ties; the final five chunks are supplied
to the answer model.

\subsubsection{Canonical Independent HyDE Retrieval}
HyDE~\cite{gao2023precise} independently generates a Spanish clinical passage
from each question using OpenAI \code{gpt-5-mini-2025-08-07}, separate from
the answer model. Each passage is generated once, cached across answer-model
conditions, and embedded with BGE-M3 for exact cosine search. Only retrieved
corpus chunks, never hypothetical text, provide answer
evidence~\cite{ji2023survey}. Cache namespaces separate dataset mode,
provider, and generator model; reuse checks \code{qa_id} and a question hash.

\subsubsection{No RAG Control}
The No RAG control returns no retrieved contexts and
omits the retrieved-context block from the answer prompt. Context-dependent
retrieval metrics are therefore not defined and are not assigned a value of
zero.

\subsection{Inference and Prompt Controls}
All answer prompts are written in Spanish. RAG prompts contain a clearly
delimited context block with positional labels and the question. Provenance
is retained in benchmark records rather than asserted as content in the
answer prompt. The No RAG condition receives the same question
formatting and chat template but omits the context block.

Inference uses the model's chat template and greedy decoding by default,
with a maximum of 2,048 newly generated tokens. Sampling is disabled unless a
run explicitly records a different configuration. QLoRA conditions load the
declared adapter on the corresponding base model; base conditions do not.
Retrieval generation and answer generation remain separate stages.

\subsection{Evaluation Framework and Metrics}
Automatic evaluation uses RAGAS~\cite{es2024ragas}.
Response-quality, retrieval-quality, and operational measurement domains
are kept separate before paired contrasts and interaction analysis.

Table~\ref{tab:results} reports six RAGAS metrics: Context Precision,
Context Recall, Faithfulness, Answer Relevancy, Answer Correctness, and
Semantic Similarity. Context Recall assesses coverage of the reference
answer by retrieved contexts; it does not require gold chunk identifiers.
Reference-chunk identity checks are separate and unavailable for the
challenge set.

The final aggregate exports do not retain the run-level evaluator
configuration, so the evaluator model, embedding model, completion limit,
and RAGAS version cannot be verified from these files.

Operational variables include input, output, and total token counts;
retrieval, generation, and end-to-end latency; and output tokens per
second. Wall-clock output rate is descriptive and not treated as
hardware-independent throughput; it divides output-token count by
generation time in seconds.

\subsection{Analysis Plan}
The primary analysis is descriptive and item-level. Means are computed from
the 10 question-level scores for each applicable metric within each balanced
scenario; context-dependent metrics are not applicable to No RAG.
Comparisons are exploratory. QLoRA is compared with
base within each retrieval strategy and family; retrieval strategies are
contrasted within adaptation and family; and model families are compared
under matched conditions.

The summary used the same 10 \code{qa_id} values shared by all 12
conditions in each scenario.

Approximate marginal 95\% Student-$t$ confidence intervals use exported
means, sample standard deviations, and counts (10 questions): the mean plus
or minus the two-sided critical value with nine degrees of freedom times
the standard error. This independent-question approximation is limited by
nonrandom selection and shared source documents.
Intervals for all metrics and unstandardized mean differences accompany the
reproducibility files. Paired tests, bootstrap intervals, and standardized
paired effects require current question-level scores, unavailable in the
aggregate exports, and are not reported.

\subsection{Reproducibility and Experimental Scope}
Reproducibility requires recording the MaternaQA-es revision,
retrieval-corpus hash, complete configuration, prompts, model and adapter
identifiers, HyDE provider settings, cache identity, evaluator version, and
decoding settings. The pinned BGE-M3 revision and deterministic RRF
tie-breaking are part of this record. Reporting follows the transparency
goals of TRIPOD-LLM and Data Statements for
NLP~\cite{gallifant2025tripod,bender2018data}.

QLoRA fine-tuning and local neural-model computations were performed on an
NVIDIA GeForce RTX 5070 Ti Laptop GPU with 12~GB of dedicated VRAM.
The comparison comprised 20 questions---10 from MaternaQA-es and 10 from
\code{sample10}---evaluated across all 12 conditions, and all principal
comparisons are made within each scenario.

% =====================================================================
\section{Results}
% =====================================================================
\subsection{Available Runs and Analysis Set}
Table~\ref{tab:results} reports the mean metrics for both scenarios;
Table~\ref{tab:correctness-ci} reports Answer Correctness confidence intervals.

% --- Table III: full-width, caption above, no vertical rules ---------
\begin{table*}[!t]
\footnotesize
\centering
\begin{threeparttable}
\caption{Mean Automatic Metrics for the 10 Shared Questions in Each Scenario}
\label{tab:results}
\begin{tabular*}{\textwidth}{@{\extracolsep{\fill}}lllrrrrrr@{}}
\toprule
\textbf{Model} & \textbf{State} & \textbf{Strategy} &
\textbf{\shortstack[r]{Answer\\Correctness}} &
\textbf{\shortstack[r]{Answer\\Relevancy}} &
\textbf{\shortstack[r]{Semantic\\Similarity}} &
\textbf{Faithfulness} &
\textbf{\shortstack[r]{Context\\Precision}} &
\textbf{\shortstack[r]{Context\\Recall}} \\
\midrule
\multicolumn{9}{c}{\textit{Panel A. MaternaQA-es test subset}} \\
\midrule
Gemma 4 & Base & No RAG & 0.372 & 0.696 & 0.704 & \dash & \dash & \dash \\
Gemma 4 & Base & Hybrid & \textbf{0.649} & 0.669 & 0.820 & \textbf{0.967} & 0.908 & 0.800 \\
Gemma 4 & Base & HyDE & 0.560 & 0.685 & 0.796 & 0.954 & \textbf{0.950} & 0.850 \\
Gemma 4 & QLoRA & No RAG & 0.384 & \textbf{0.776} & 0.726 & \dash & \dash & \dash \\
Gemma 4 & QLoRA & Hybrid & 0.605 & 0.759 & \textbf{0.838} & 0.883 & 0.908 & 0.800 \\
Gemma 4 & QLoRA & HyDE & 0.637 & 0.639 & 0.818 & 0.882 & \textbf{0.950} & 0.750 \\
MedGemma 1.5 & Base & No RAG & 0.388 & 0.611 & 0.714 & \dash & \dash & \dash \\
MedGemma 1.5 & Base & Hybrid & 0.599 & 0.648 & 0.794 & 0.897 & \textbf{0.950} & \textbf{0.900} \\
MedGemma 1.5 & Base & HyDE & 0.595 & 0.737 & 0.807 & 0.842 & \textbf{0.950} & \textbf{0.900} \\
MedGemma 1.5 & QLoRA & No RAG & 0.330 & 0.683 & 0.720 & \dash & \dash & \dash \\
MedGemma 1.5 & QLoRA & Hybrid & 0.609 & 0.700 & 0.826 & 0.883 & 0.908 & 0.850 \\
MedGemma 1.5 & QLoRA & HyDE & 0.574 & 0.674 & 0.828 & 0.859 & \textbf{0.950} & \textbf{0.900} \\
\midrule
\multicolumn{9}{c}{\textit{Panel B. Obstetric challenge set}} \\
\midrule
Gemma 4 & Base & No RAG & 0.359 & 0.716 & 0.651 & \dash & \dash & \dash \\
Gemma 4 & Base & Hybrid & \textbf{0.457} & 0.554 & 0.689 & \textbf{0.966} & 0.800 & 0.400 \\
Gemma 4 & Base & HyDE & 0.335 & 0.658 & \textbf{0.707} & 0.902 & 0.958 & 0.200 \\
Gemma 4 & QLoRA & No RAG & 0.205 & 0.845 & 0.634 & \dash & \dash & \dash \\
Gemma 4 & QLoRA & Hybrid & 0.349 & 0.736 & 0.693 & 0.838 & 0.775 & 0.400 \\
Gemma 4 & QLoRA & HyDE & 0.399 & 0.672 & \textbf{0.707} & 0.919 & 0.933 & 0.200 \\
MedGemma 1.5 & Base & No RAG & 0.287 & 0.832 & 0.631 & \dash & \dash & \dash \\
MedGemma 1.5 & Base & Hybrid & 0.402 & 0.368 & 0.622 & 0.850 & 0.842 & \textbf{0.500} \\
MedGemma 1.5 & Base & HyDE & 0.260 & 0.545 & 0.639 & 0.850 & \textbf{1.000} & 0.200 \\
MedGemma 1.5 & QLoRA & No RAG & 0.234 & \textbf{0.913} & 0.649 & \dash & \dash & \dash \\
MedGemma 1.5 & QLoRA & Hybrid & 0.455 & 0.787 & 0.678 & 0.732 & 0.842 & 0.400 \\
MedGemma 1.5 & QLoRA & HyDE & 0.316 & 0.725 & 0.668 & 0.771 & \textbf{1.000} & 0.200 \\
\bottomrule
\end{tabular*}
\begin{tablenotes}[flushleft]
\item[] \textit{Note.} Means are calculated from 10 responses per condition
and scenario for each applicable metric. Dashes indicate metrics not applicable
without retrieval. Bold values identify each metric's maximum within a scenario.
\end{tablenotes}
\end{threeparttable}
\end{table*}

The final analysis comprised 240 condition-item outputs: 120 for the
MaternaQA-es subset and 120 for the obstetric challenge set. Each scenario
contained 10 questions evaluated across all 12 conditions, yielding 10 final
responses per condition. All automatic metrics applicable to each condition
were available; no scores were imputed.

\subsection{MaternaQA-es Subset}
Gemma~4 Base with Hybrid retrieval achieved the highest observed mean Answer
Correctness (0.649), whereas Gemma~4 QLoRA with No RAG achieved the highest
Answer Relevancy (0.776), and Gemma~4 QLoRA with Hybrid retrieval achieved the
highest Semantic Similarity (0.838). Under No RAG, QLoRA changed Gemma~4
Answer Correctness, Answer Relevancy, and Semantic Similarity by +0.012,
+0.080, and +0.022, respectively. For MedGemma~1.5, the corresponding changes
were $-$0.058, +0.072, and +0.006.

Relative to No RAG, Hybrid and HyDE increased Answer Correctness and Semantic
Similarity in all four model-state combinations. Answer Relevancy increased
in two of four combinations with Hybrid and one of four with HyDE. Compared
directly with Hybrid, HyDE increased Answer Correctness in one of four
combinations, Answer Relevancy in two, and Semantic Similarity in two,
indicating heterogeneous effects across models and metrics.

\subsection{Obstetric Challenge Set}
Within the evaluated obstetric challenge subset, Gemma~4 Base with Hybrid
retrieval achieved the highest observed mean Answer Correctness (0.457), and
MedGemma~1.5 QLoRA with No RAG achieved the highest Answer Relevancy (0.913).
Gemma~4 Base and Gemma~4 QLoRA with HyDE tied for the highest Semantic
Similarity (0.707). Under No RAG, QLoRA changed Answer Correctness, Answer
Relevancy, and Semantic Similarity by $-$0.154, +0.129, and $-$0.017 for
Gemma~4, and by $-$0.054, +0.081, and +0.017 for MedGemma~1.5.

Relative to No RAG, Hybrid increased Answer Correctness in all four model
conditions but reduced Answer Relevancy in every condition. Semantic
Similarity increased for both Gemma~4 conditions and MedGemma~1.5 QLoRA, but
decreased for MedGemma~1.5 Base.

Relative to No RAG, HyDE increased Answer Correctness for Gemma~4 QLoRA
(+0.194) and MedGemma~1.5 QLoRA (+0.082), but reduced it for Gemma~4 Base
($-$0.024) and MedGemma~1.5 Base ($-$0.028). Answer Relevancy decreased in
all four conditions, while Semantic Similarity increased in all four.

Compared directly with Hybrid retrieval, HyDE improved Answer Correctness
for Gemma~4 QLoRA (+0.050) and Answer Relevancy for Gemma~4 Base (+0.104)
and MedGemma~1.5 Base (+0.177). Semantic Similarity increased in three of
four model conditions. These descriptive patterns show that higher retrieval
scores did not guarantee uniformly better response metrics.

% --- Table IV: full-width, caption above, no vertical rules ---------
\begin{table*}[t]
\footnotesize
\centering
\begin{threeparttable}
\caption{Answer Correctness: Mean and Approximate Marginal 95\% Confidence
Interval}
\label{tab:correctness-ci}
\begin{tabular*}{\textwidth}{@{\extracolsep{\fill}}lccc@{}}
\toprule
\textbf{Model and state} & \textbf{No RAG} & \textbf{Hybrid} & \textbf{HyDE} \\
\midrule
\multicolumn{4}{c}{\textit{MaternaQA-es test subset}} \\
\midrule
Gemma 4 Base        & 0.372 [0.238, 0.506] & 0.649 [0.456, 0.841] & 0.560 [0.375, 0.744] \\
Gemma 4 QLoRA       & 0.384 [0.269, 0.500] & 0.605 [0.422, 0.788] & 0.637 [0.441, 0.833] \\
MedGemma 1.5 Base   & 0.388 [0.258, 0.517] & 0.599 [0.398, 0.801] & 0.595 [0.387, 0.803] \\
MedGemma 1.5 QLoRA  & 0.330 [0.213, 0.446] & 0.609 [0.454, 0.764] & 0.574 [0.401, 0.747] \\
\midrule
\multicolumn{4}{c}{\textit{Obstetric challenge set}} \\
\midrule
Gemma 4 Base        & 0.359 [0.159, 0.558] & 0.457 [0.244, 0.671] & 0.335 [0.125, 0.545] \\
Gemma 4 QLoRA       & 0.205 [0.127, 0.282] & 0.349 [0.185, 0.513] & 0.399 [0.254, 0.544] \\
MedGemma 1.5 Base   & 0.287 [0.164, 0.411] & 0.402 [0.139, 0.664] & 0.260 [0.123, 0.396] \\
MedGemma 1.5 QLoRA  & 0.234 [0.152, 0.315] & 0.455 [0.249, 0.660] & 0.316 [0.098, 0.533] \\
\bottomrule
\end{tabular*}
\begin{tablenotes}[flushleft]
\item[] \textit{Note.} Mean [approximate marginal 95\% confidence interval];
$n = 10$ questions. Student-$t$ intervals
assume independent questions; they do not test paired differences.
\end{tablenotes}
\end{threeparttable}
\end{table*}

\subsection{Retrieval-Grounding Metrics}
In the MaternaQA-es subset, Hybrid Context Precision ranged from 0.908 to
0.950, Context Recall from 0.800 to 0.900, and Faithfulness from 0.883 to
0.967. HyDE Context Precision was 0.950 in all four model states, Context
Recall ranged from 0.750 to 0.900, and Faithfulness from 0.842 to 0.954.
HyDE exceeded the corresponding Hybrid Context Precision in three of four
model states and increased Context Recall in two of four.

In the obstetric challenge set, Hybrid Context Precision ranged from 0.775
to 0.842, Context Recall from 0.400 to 0.500, and Faithfulness from 0.732
to 0.966. HyDE Context Precision ranged from 0.933 to 1.000, Context Recall
was 0.200 in all four model states, and Faithfulness ranged from 0.771 to
0.919. HyDE achieved higher Context Precision than Hybrid for all four model
states, while Context Recall was lower for all four.

Across both scenarios, the 20 unique uncached HyDE generations averaged
18.5 seconds, 67.1 input tokens, and 1,437.3 API output tokens. The same
hypothetical documents were then reused from cache across the remaining
answer-model conditions, reducing model-dependent variation in the
query-expansion stage.

\subsection{Uncertainty of Answer Correctness}
Table~\ref{tab:correctness-ci} reports broad approximate marginal 95\%
confidence intervals for Answer Correctness, indicating substantial
item-level variability. They do not quantify paired differences, and overlap
is not a significance test. Rankings remain exploratory; statistically
significant superiority is not established.

% =====================================================================
\section{Discussion}
Across the evaluated 20-question benchmark, the observed effect of QLoRA
under No RAG varied by model family and scenario. Answer Relevancy increased
in all four model-scenario comparisons. Answer Correctness decreased for both
model families in the obstetric challenge set and for MedGemma~1.5 in the
MaternaQA-es subset; Semantic Similarity decreased for Gemma~4 in the challenge
set but increased for MedGemma~1.5 and for both models in MaternaQA-es.

Within the evaluated questions, medical pretraining did not show a uniform
advantage after domain-specific adaptation across the tested metrics and
configurations. These findings support evaluating multiple
dimensions rather than treating a single automatic score as a complete
measure of quality.

Hybrid retrieval increased Answer Correctness in all four model conditions
in each scenario but reduced Answer Relevancy in six of eight comparisons.
One plausible interpretation is that retrieved evidence increased agreement
with reference answers while sometimes reducing direct focus. Retrieval-
grounding scores were generally high for Context Precision but more variable
for Context Recall, indicating that relevant evidence was often surfaced
without consistently covering all reference-supporting information.

HyDE's high RAGAS Context Precision did not translate into consistent
improvements in final answers. Within the evaluated questions, retrieval
precision alone is therefore insufficient to favor HyDE over Hybrid;
answer-quality metrics must also be considered.

% =====================================================================
\section{Conclusions}
For the models and conditions evaluated, MedGemma did not consistently
outperform Gemma after grounded QLoRA adaptation. The study therefore does not
establish a general advantage for the medically pretrained model once both
models receive the same domain-specific supervision. QLoRA also showed no
consistent benefit in Answer Correctness, despite improving Answer Relevancy
under No RAG.

The main implication is that medical specialization and domain adaptation
cannot substitute for evaluating the complete question-answering
configuration. Hybrid retrieval produced more consistent correctness gains
than QLoRA in this benchmark, but neither adaptation nor retrieval improved
every assessed outcome. Configuration choices therefore require evidence from
the intended task rather than assumptions based on medical specialization or
fine-tuning alone.

These conclusions remain exploratory: they concern two model families, ten
questions per scenario, and automated metrics, with relevant source documents
available for retrieval. The comparison does not isolate the causal effect of
medical pretraining, demonstrate model equivalence, or establish clinical
reliability. Broader paired evaluation and clinician assessment are needed to
determine whether these findings generalize.
% =====================================================================
\section{Limitations and Future Work}
The approximate marginal intervals do not establish paired statistical
significance and do not account for document clustering, subset-selection
bias, or evaluator variability.

% =====================================================================
%  Bibliography -- IEEEtran.bst, numeric citation order
% =====================================================================
\bibliographystyle{IEEEtran}
\bibliography{references}

\end{document}
